\section{Related work}
\label{sec:rw}

\subsection{Capability operating systems}

Capabilities trace to Dennis and Van~Horn~\cite{dennis1966}, with
Hydra and StarOS carrying the idea through the 1970s. Redell's
thesis~\cite{redell74} contributed the epoch-versioned credential,
the ancestor of every revocable-capability scheme since, ours
included. KeyKOS, EROS~\cite{shapiro1999eros}, and
CapROS~\cite{capros} ran the pattern at OS scale, Coyotos refined the
typing discipline, and seL4~\cite{klein2009sel4} finished the arc
with full functional-correctness verification.

More recent work splits on whether capability identifiers should be
\textit{names} (KeyKOS-style) or \textit{addresses} (CHERI-style).
Our tokens are names that smuggle in an address-like hint---the slot
index of \Cref{sec:design:fastpath}---which buys the cache behaviour
of addresses without surrendering the unforgeability of names.

\subsection{Hardware capability systems}

CHERI~\cite{cheri-asplos15,cheri-tcg22} is the most developed
hardware-capability platform, and CAP-VMs~\cite{capvms} builds
capability isolation for cloud microservices on top of
CHERI/Morello. Cornucopia\,Reloaded~\cite{cornucopia-reloaded} is the
modern peer for the revocation pattern itself; we use their
algorithm, realised in software on stock Linux.
PICASSO~\cite{picasso} (S\&P~'24) reaches a similar bulk-invalidation
effect through hardware provenance.

\subsection{Userspace IPC primitives}

Aeron~\cite{aeron-real-logic} publishes the best single-host IPC
latency we are aware of, and ByteDance's
shmipc~\cite{bytedance-shmipc} deploys the same architectural idea in
production service meshes. io\_uring~\cite{iouring} put the
template---mmap'd SQ/CQ rings, futex-style notification---inside the
kernel; our ring buffer reuses it above the kernel. eRPC~\cite{erpc}
(NSDI~'19) holds the lowest peer-reviewed RPC latency we know
($\sim$2.3\,$\mu$s single-hop over DPDK).
LITESHIELD~\cite{liteshield} (USENIX ATC~'25) is the most recent
userspace microkernel-style runtime on Linux; the structural overlap
with \sysname is real, and the capability layer is what we add on
top of that shape. Cap'n~Proto's RPC layer~\cite{capnproto-rpc}
lives higher in the stack---marshalling and authorisation over a
transport---so we do not treat it as a raw-RTT comparison.

\subsection{Kernel-resident capability systems on Linux}

HongMeng~\cite{hongmeng} (OSDI~'24) and CertiKOS~\cite{certikos} are
the closest kernel-resident capability systems with Linux-compatible
ambitions. We sit above them in the stack rather than competing
inside it, and our framing leans on HongMeng's own account of the
trade-off between pure capability chains and address tokens
(\Cref{sec:disc:hongmeng}).

\subsection{Confidential-compute capability binding}

Confidential-compute work has started tying capability mediation to
attestation: ScaleTrust~\cite{scaletrust} explores TDX-bound
credentials, Conclave~\cite{conclave} SEV-SNP-sealed objects. We
deliberately leave that binding to follow-up work (M-23 ConfCompute)
and keep this paper inside the Linux process boundary.

\subsection{LLM-agent sandboxing}

Anthropic's Computer Use~\cite{anthropic-cu}, the OpenAI Agents
SDK~\cite{openai-agents}, and Devin~\cite{devin} all sandbox agent
tool calls today, and none exposes a capability layer above the
syscall boundary. With NIST AI~600-1~\cite{nist-ai-600-1} and EU AI
Act Article~15~\cite{euaiact15} demanding verifiable mediation, we
read \sysname as the substrate that demand will eventually pull into
existence; the AgentGuard module (M-22) is in preparation as a
follow-up paper.

\subsection{Formal verification}

We use CBMC~\cite{cbmc} for the monotonicity proof.
SPIN~\cite{spin} could have handled an LTL formulation, but our
property is a quantifier-free boolean over bitmasks and CBMC states
it directly. The scope---one property, a 280-line C model, 37
checks---is orders of magnitude below the seL4
effort~\cite{klein2009sel4}, and we call it what it is:
spec correspondence, not functional correctness.
